\subsection{Retained structural-initialization evidence}\label{sec:mp_legacy_frontend}

The frozen combined legacy frontend has a completed component study on
FactorShock, InventorySupplyChain, and QueueResourceControl.  At nominal
dimension 1,000, each family contributes 20 algorithmic seeds.  The frontend
uses 384 reusable source calls per held-out-family archive to freeze ten
target proposals.  Proposal-only selection uses ten target search calls;
Stacked GP and canonical SAASBO add three adaptive calls.  All three frontend
arms achieved 60/60 true-feasible and independently certified recommendations
with zero false certificates.  Replacing the frontend by common Sobol gave
0/60, 1/60, and 0/60 certifications, respectively
(Table~\ref{tab:mp_legacy_frontend}).
This frontend combines universal policy shapes with source templates through
the historical \texttt{LearnedMetaPrior} implementation; it is distinct from
the later public 64-profile \texttt{SourceScoredProfileAtlas}.  Component
ablations certified 40/60 with source templates alone (0/20 FactorShock,
20/20 Inventory, 20/20 Queue), and 23/60 with universal shapes alone
(20/20, 2/20, 1/20).  The combined result therefore reflects complementary
components rather than a contribution attributable wholly to source learning.

\begin{table}[htbp]
\centering
\small
\caption{Historical frontend component study. Source calls are per reusable
archive and search calls are per target run; terminal verification calls are
additional and separately recorded. These results use the common historical
v69 Gaussian terminal protocol, not the current physical-Inventory verifier.}
\label{tab:mp_legacy_frontend}
\resizebox{\textwidth}{!}{\begin{tabular}{llrrrr}
\toprule
Frontend & Backend & Source & Search & Certified & False \\
\midrule
Combined legacy frontend & Proposal only & 384 & 10 & 60/60 & 0 \\
Combined legacy frontend & Stacked GP & 384 & 13 & 60/60 & 0 \\
Combined legacy frontend & Canonical SAASBO & 384 & 13 & 60/60 & 0 \\
Common Sobol & Proposal only & 0 & 10 & 0/60 & 0 \\
Common Sobol & Stacked GP & 384 & 13 & 1/60 & 0 \\
Common Sobol & Canonical SAASBO & 0 & 13 & 0/60 & 0 \\
\bottomrule
\end{tabular}
}
\end{table}

The native transfer matrix retains each comparator's own initialization and
adaptation with the 384-call source archive and 13-call target search budget.
Its best certification totals were RGPE's 30/60 and MetaBO's 28/60.  In the
separate common-frontend backend comparison, the eight transfer methods reached
56--60/60.  These comparisons retain both the native and supplied-frontend
pipelines; the common-Sobol factorial above supplies the direct frontend
intervention for the two fitted backends.  The supplement preserves the
individual method totals and evidence locations.

This study supports the value of structural initialization on these three
fixed task families.  Its seeds measure algorithmic repeatability, not
generalization across 60 independent task laws.  The current physical
Inventory model has four state-bin actions and replicated mean/risk labels;
it uses a common eight-policy initialization rather than the combined frontend.
Consequently, the frontend remains a demonstrated component of the project,
while independent decision superiority of the combined frontend--function-prior--KG
pipeline remains unestablished.  A development integration check reuses the
historical frozen proposals and source archives at seed 80 in each family,
with ten common single-episode observations and three subsequent raw-observation
KG calls.  The source-orthogonal and same-source full priors do not improve
the deployed objective over proposal-only selection.  Both can recommend
infeasible unobserved Inventory candidates after ten calls; a common independent
verifier retaining initial-policy support preserves the deployed result.
Under the declared primary-certificate utility and common finite candidate
bank, the retained initial deployment leaves no utility headroom in FactorShock
and only 0.0021 and 0.0014 in Inventory and Queue, respectively.  These
post-decision bounds apply to this candidate bank and proxy utility.
This check uses a declared primary-certificate utility for acquisition and a
guarded shortlist for deployment; it diagnoses the remaining mismatch between
these objectives, rather than establishing a unified deployment KG advantage.
Its replay entry point is given in the supplement.
